\chapter{动机与行为权限边界}

\section{动机是面向未来状态的认知}

动机（Motivation）的当前目标范围包括愿望（Desire）、目标（Goal）、关切（Concern）和优先级（Priority）：主体偏好什么未来状态、已形成什么目标、持续关注什么问题，以及多个未来状态如何排序。这些属于认知状态，需要有来源、修订和生命周期；具体字段及满足、放弃、冲突语义仍待确定。\cite{nous-research-package}

\section{主体目标与行为承诺}

Nous Wave 动机领域中的目标表示主体期待的未来状态。有约束力的承诺或义务属于 Heptalogos 的行为和规范状态；执行以及对外效果也由 Heptalogos 的相应领域管理。

\begin{center}
\small
主体目标修订 $\rightarrow$ Heptalogos 审查与采纳 $\rightarrow$
有约束力的承诺 $\rightarrow$ 执行结果回执 $\rightarrow$
观察事件或使用事件
\end{center}

工具调用、排程和对外消息由 Heptalogos 的行为授权合同管理。信念—愿望—意图（BDI）研究区分信念、愿望和意图；社会承诺研究进一步讨论承诺的创建、满足和取消，为动机与行为权限划分提供概念背景。\cite{rao-georgeff1995,singh1997,singh1999}

\section{行为与认知权限}

\begin{ReportWideTable}{>{\raggedright\arraybackslash}p{5.0cm}>{\raggedright\arraybackslash}p{8.5cm}}
\captionof{table}{Nous Wave 动机认知与 Heptalogos 行为权限的边界}\label{tab:goal-commitment}
\toprule
\textbf{语义项} & \textbf{归属与效果}\\
\midrule
目标、愿望、关切 & 属于 Nous Wave 动机认知，可检索、修订或转为提案。\\
承诺、义务 & 属于 Heptalogos 行为或规范状态，依该系统的审查与生命周期规则处理。\\
通信承诺（CommunicationCommit） & Heptalogos 对本地主体聊天回复的批准记录。\\
外部效果 & 由获授权的执行边界处理，不从记忆、自我认知或自然语言直接产生。\\
\bottomrule
\end{ReportWideTable}

\section{开放范围}

动机领域的最小职责已进入目标集合。后续合同需确定价值观归属、需要与兴趣的语义、目标满足或放弃状态，以及优先级作为持久认知还是运行时计算的具体方式。\cite{nous-research-package}
